% master.tex : master-fil for projektet
% ------------------------------------------------------------------------------
% Dette er hovedfilen for projektet, hvori indhold fra alle input-filer (tekst,
% billeder, litteraturdatabaser, osv.) samles

% Dokumenttypen 'book' er valgt pga. dens mange fleksible indstillinger
% Se https://tex.stackexchange.com/a/36989/118167
\documentclass[11pt,a4paper,twoside,openright,danish]{book}

% Variabler, som bruges til automatisk at indsætte titel, forfattere, osv. på
% forsiden og titelbladet.
\def \projecttitle       {P1-projekt 2020 for matematik-økonomi}
\def \projectsubtitle    {Optimering af gaslager}
\def \projecttheme       {Optimering af gaslager}
\def \projectdegree      {Matematik-Økonomi}  % eller Matematik-Økonomi/Teknologi
\def \projectperiod      {Efterårssemesteret 2020}
\def \projectnumber      {P1}
\def \projectgroup       {B350}
\def \projectauthors     {
  Beate Marcussen\\
  Marcus Søndergaard Aaskoven\\
  Mathias Graversen\\
  Mattias Adam Arvidsson\\
  Vennan Vithiyatharan
  % ...
}
\def \projectsupervisors {
  PJanus Valberg-Madsen\\
  % ...
}

% Preamblet indeholder alle de indstillinger og makroer, som skal indsættes for
% hovedindholdet, og i denne skabelon samles det i filen aaumath.sty, som
% definerer en pakke, der kan indlæses med \usepackage.
\usepackage{aaumath}

% Dokumentets indhold indsættes mellem \begin- og \end-makroerne for
% 'document'-blokken
\begin{document}

% Dokumentets 'front matter' tælles ikke med ifm. antal sider og nummereres med
% romerske tal. Herunder hører f.eks. forsiden, titelbladet, forordet og
% indholdsfortegnelsen.
\frontmatter
% incl/misc/frontpage.tex : rapportens forside
% ------------------------------------------------------------------------------


\backgroundsetup{
  scale = 1,
  angle=0,
  opacity=1,
  contents = {
    \includegraphics[width=\paperwidth,height=\paperheight]{fig/img/aau/waves.pdf}
  }
}
\BgThispage
\pdfbookmark[0]{Forside}{forside}
\begin{titlepage}
  \centering
  \phantom{}
  \vspace{2cm}

  % AAU-segl
  \begin{minipage}[c]{0.2\paperwidth}
    \centering
    \makebox[0pt]{
      % fig/tikz/aau-badge.tex : AAU-logo til forsiden
% ------------------------------------------------------------------------------

\begin{tikzpicture}
  % Tegn hvid cirkel og tilføj det gennemsigtige, blå logo ovenpå
  \node[circle,color=white,fill=white,minimum size=1.175\textwidth] at (0,0) {};
  \node at (0,0) {\includegraphics[width=\textwidth]{fig/img/aau/logo-circle.pdf}};
\end{tikzpicture}

    }
  \end{minipage}

  % Hovedindhold
  \vspace{4cm}
  {\fontfamily{bch}\selectfont
    \fboxsep0pt\colorbox{white}{
      \begin{minipage}{\textwidth}
        \centering
        \color{AAUblue1}

        \vspace{2em}
        {\Huge\bfseries\projecttitle}

        {\Large\bfseries\projectsubtitle}

        \bigskip
        \parbox{\textwidth}{\centering\large\projectauthors}

        \bigskip
        {\bfseries\large{\projectnumber}-Projekt, Gruppe \projectgroup, \projectdegree}
        \vspace{2em}
      \end{minipage}
    }
  }

\end{titlepage}

% incl/misc/titlepage.tex : rapportens titelblad
% ------------------------------------------------------------------------------
% Titelbladet genereres af makroen \aautitlepage, som er defineret i
% /incl/pre/ext/aautitlepage.sty


\pdfbookmark[0]{Titelblad}{titelblad}
\aautitlepage{
  \projectinfo{
    \projecttitle
  }{
    \projecttheme
  }{
    \projectperiod
  }{
    Gruppe \projectgroup
  }{
    \parbox[t]{\textwidth}{\projectauthors}
  }{
    \parbox[t]{\textwidth}{\projectsupervisors}
  }{
    \today
  }
}{
  \textbf{Institut for Matematiske Fag}\\
  Skjernvej 4A\\
  DK-9220 Aalborg Ø\\
  \href{http://math.aau.dk}{http://math.aau.dk}
}{
  % incl/misc/abstract.tex : projektets abstract 

This study investigates how algorithms and graph theory can help with the optimization of profit from a gas storage, over the course of a given period of time, by buying and selling gas units. 

In order to give clarification of the calculations and methods used in this study, the essential and relevant graph theory, and theory of algorithms will therefore be explained.
During this study it was found that Dijkstra's algorithm was very useful for solving shortest path problems in graphs. This study implements Dijkstra's algorithm in Python 3 in order to find the longest path through a graph. This can be done by finding the shortest path through the inverted non-negative graph. This path will be the path that yields the greatest profit. In addition, this path is the one that will give us the most optimal trading strategy.

The greatest profit in the original problem is 252.73€, and in the extended problem, it is 188.27€. From this, it can be concluded that the extensions that have been made, significantly reduce the profit. This is mainly due to the fluctuating limits of inventory, and the limits of buying and selling per month, as they affect the profit over the course of several months, compared to the penalty factor which only affects the last two months.

}

\tableofcontents

% Dokumentets 'main matter' (hovedindhold) er der, hvor det meste indhold skal
% sættes ind. Sider og overskrifter er nummererede med arabiske tal.
\mainmatter
\chapter{Forord}
Følgende projekt er udarbejdet af gruppe B350 bestående af fem stud.scient.oeconer på 1. semester. Projektet er skrevet i efteråret 2020 og beskæftiger sig med diskret matematik herunder optimering af et gaslager som basisproblem. Diskret matematik indgår i studiets læreplan og er derfor et relevant emne til projektskrivning. Projektet er skrevet i \LaTeX, og delt via GitAhead.
Vi vil som gruppe rette en stor tak til vores vejleder Janus Valberg-Madsen for godt samarbejde gennem projektforløbet samt gode råd og rettelser.


\chapter{Indledning}


\chapter{Redegørelse}
\section{Sektion 1}

\section{Sektion 2}

\section{Sektion 3}

\chapter{Analyse}
\section{Sektion 1}

\section{Sektion 2}

\section{Sektion 3}

\chapter{Vurdering af metoderne}
\section{Sektion 1}

\section{Sektion 2}

\section{Sektion 3}

\chapter{Konklusion}
% Input-filer bør opdeles således, at hver fil svarer til et kapitel. Makroen
% \include indsætter et sideskift og indholdet fra den givne stil.

% \include{incl/main/example1}
% \include{incl/main/example2}
% ...

% Appendicer indsættes inde i en appendices-blok og bliver nummereret med
% bogstaver i stedet for tal
\begin{appendices}
  % \include{incl/app/appendix1}
  % \include{incl/app/appendix2}
  % ..
\end{appendices}

% Dokumentets 'back matter' er til ekstra ting som f.eks. litteraturlisten.
% Overskrifter bliver ikke nummereret her.
\backmatter

% Automatisk litteraturliste baseret på, hvilke kilder, der er blevet refereret
% til i løbet af rapporten.
\bibliographystyle{apalike}
\bibliography{
  incl/bib/books,
  incl/bib/articles,
  incl/bib/software
}

\end{document}
